\tema{e) El código: qué archivo abrir y cómo ejecutarlo}
\paso{1. Ya no se necesita buscar un archivo que no está adjunto.}
El programa didáctico completo está en
\begin{center}\texttt{ejemplos/paso\_a\_paso.py}\end{center}
y se reproduce íntegro al final de este documento. Incluye importaciones, parámetros, todas las funciones matemáticas, cálculo de métricas, escritura de tablas y generación de cuatro imágenes. No es solo el núcleo de unas veinte líneas. No depende de importar las funciones internas del paquete \texttt{feal}.

\paso{2. Preparar el entorno.}
Descargar o clonar el repositorio y abrir una terminal en su carpeta. \texttt{uv} administra el intérprete de Python y sus bibliotecas. \texttt{uv.lock} fija las versiones. La primera ejecución necesita conexión para descargar lo que falte. Ejecutar:
\begin{verbatim}
uv sync --locked
uv run --locked python ejemplos/paso_a_paso.py
\end{verbatim}
La primera línea instala el entorno declarado sin modificar el archivo de versiones. La segunda ejecuta el archivo con ese entorno. \texttt{python} es el intérprete; la ruta que le sigue es el programa que debe leer.

\paso{3. Encontrar los resultados.}
Aparecen en \texttt{output/didactico/}: \texttt{metricas.csv}, \texttt{coeficientes\_gamma.csv}, \texttt{actividad.png}, \texttt{errores.png}, \texttt{termodinamica.png} y \texttt{serie\_infinita.png}. La terminal también muestra las métricas. El programa guarda las figuras; no necesita abrir ventanas.

\paso{4. Regenerar el documento y los desgloses completos.}
El comando integral es:
\begin{verbatim}
uv run --locked feal
\end{verbatim}
Este comando usa los módulos de \texttt{src/feal/}, genera tablas y figuras con formato editorial y compila LaTeX. Requiere \texttt{pdflatex}; el README indica cómo instalarlo. Con \texttt{uv run --locked feal --solo-datos} se omite la compilación. El programa didáctico no necesita LaTeX.

\paso{5. Python y MATLAB son lenguajes diferentes.}
El archivo entregado es Python. Pegar su contenido en MATLAB no lo ejecuta. Para trasladar el método, hay que traducir sintaxis: Python cuenta índices desde 0 y MATLAB desde 1; NumPy aplica \texttt{x**k} elemento a elemento, mientras que en MATLAB se usa \texttt{x.\^{}k}. Los CSV sí pueden abrirse desde MATLAB. Esta edición incluye también una versión completa en MATLAB, explicada y listada en las secciones posteriores. Su compatibilidad se comprueba con GNU Octave.

\tema{e) Leer Python desde cero: signos, arreglos e índices}
\paso{1. Asignación y comparación.}
En código, \texttt{x = 0.1} guarda el número 0.1 con el nombre \texttt{x}. No es una ecuación que deba resolver el programa. \texttt{==} compara igualdad; \texttt{<} compara menor que y \texttt{<=} menor o igual. El resultado de una comparación es verdadero o falso. \texttt{**} eleva a una potencia y \texttt{*} multiplica.

\paso{2. Comentarios y nombres.}
Lo que aparece después de \texttt{\#} explica la línea y no se ejecuta. Los nombres \texttt{np} y \texttt{plt} son abreviaturas definidas en las importaciones: representan NumPy y la interfaz de gráficas. Las cadenas entre comillas son texto, no nombres de variables. \texttt{f"Orden \{orden\}"} inserta el valor de \texttt{orden} en ese texto.

\paso{3. Un arreglo almacena varios números.}
\texttt{C = np.array([-3.8685, 4.773, 2.06, -2.964])} crea cuatro posiciones. \texttt{C[0]} vale $-3.8685$; \texttt{C[1]} vale $4.773$; \texttt{C[2]} vale $2.06$; \texttt{C[3]} vale $-2.964$. La primera posición es 0. \texttt{len(C)} sería 4.

\paso{4. Operaciones sobre arreglos.}
Si \texttt{x} contiene $[0,0.1,0.2]$, su cuadrado por componentes es $[0,0.01,0.04]$. Entonces \texttt{2.06*x**2} produce $[0,0.0206,0.0824]$. Son tres cálculos independientes, uno por composición, realizados por una sola instrucción. Esta es la razón de que pocas líneas puedan generar muchos puntos.

\paso{5. Matrices e índices dobles.}
\texttt{p} tiene 4 filas y 10001 columnas. \texttt{p[2]} es toda la fila del orden 2. \texttt{p[2, 0]} sería su primer valor, en $x=0$. En general, \texttt{p[N,j]} corresponde a $P_N(x_j)$. \texttt{np.zeros((4, len(x)))} reserva esa matriz llena de ceros; no calcula todavía Taylor.

\paso{6. Funciones, ciclos y sangría.}
\texttt{def} define una función y \texttt{return} devuelve su resultado. \texttt{for} repite un bloque y \texttt{if} elige qué bloque ejecutar. Los espacios al inicio de cada línea indican a qué bloque pertenece: copiarlos mal puede cambiar el programa o causar un error. Definir una función no la ejecuta; hay que llamarla por su nombre y darle los argumentos.

\tema{e) Seguir un ciclo a mano y entender las funciones matemáticas}
\paso{1. Seguir \texttt{calcular\_logaritmo} con $x=0.1$ y orden 2.}
En el programa, $x$ es un arreglo, pero podemos seguir una sola de sus entradas. Inicialmente \texttt{resultado=0}. \texttt{range(orden+1)} es \texttt{range(3)}, que produce $k=0,1,2$; no incluye 3.
\begin{center}\begin{tabular}{rrrrr}\toprule
$k$ & Potencia & Coeficiente & Término & Resultado acumulado\\\midrule
0 & 1 & $-3.8685$ & $-3.8685$ & $-3.8685$\\
1 & 0.1 & 4.773 & 0.4773 & $-3.3912$\\
2 & 0.01 & 2.06 & 0.0206 & $-3.3706$\\\bottomrule
\end{tabular}\end{center}
Al terminar, \texttt{return resultado} entrega $-3.3706$ para esa entrada. Las otras 10000 entradas recorren las mismas operaciones.

\paso{2. Seguir \texttt{calcular\_diferencia}.}
Para orden 1, \texttt{range(orden+1,4)} produce $k=2,3$. Se inicia $\Delta=0$; primero se resta $2.06x^2$, luego se resta $-2.964x^3$. El segundo signo negativo produce una suma: $\Delta=-2.06x^2+2.964x^3$. Para orden 3, \texttt{range(4,4)} está vacío y el resultado sigue en cero.

\paso{3. Seguir \texttt{calcular\_hierro}.}
El ciclo exterior selecciona $k=1,2,3$ según el orden. El interior construye el corchete $\ln(1-x)+\sum_{j=1}^kx^j/j$. \texttt{pequeno = x < 0.1} es una máscara: una lista de verdaderos y falsos. \texttt{x[pequeno]} selecciona solo los puntos verdaderos. En esos lugares se reemplaza el corchete por la cola estable explicada antes. Finalmente se multiplica por $kc_k$ y se suma al hierro acumulado. El programa didáctico espera que $x$ sea un arreglo, como la malla que crea; no se diseñó como interfaz para escalares arbitrarios.

\paso{4. Seguir \texttt{calcular\_coeficientes\_gamma} para $n=2$.}
$\min(3,2)=2$, así que \texttt{range(1,3)} produce $k=1,2$. Para $k=1$, el producto es $1c_1b_{2-1}=c_1b_1$. Para $k=2$, es $2c_2b_{2-2}=2c_2b_0$. Se suman y dividen entre $n=2$. $b_0$ y $b_1$ ya existen porque se calcularon antes. Nunca se busca $b_{-1}$.

\tema{e) Del cálculo a los archivos y a cada gráfica}
\paso{1. Organizar resultados.}
\texttt{calcular()} reúne las funciones anteriores. \texttt{metricas.append(...)} añade una fila por orden. El diccionario devuelto asocia nombres de texto con arreglos: \texttt{resultados["gamma"]} recupera los coeficientes y \texttt{resultados["g"]} las energías. En el diccionario \texttt{q}, la clave 5 corresponde a todo el polinomio $Q_5$ evaluado en la malla.

\paso{2. Guardar una tabla.}
\texttt{guardar\_tabla} abre el archivo en modo \texttt{"w"}, que lo crea o sustituye en la carpeta de resultados. Escribe primero la cabecera y después las filas. \texttt{with} asegura que el archivo se cierre. \texttt{zip(range(13), b)} empareja el índice 0 con $b_0$, el 1 con $b_1$, y así hasta 12.

\paso{3. Dibujar dos paneles.}
\texttt{plt.subplots(1,2,...)} crea una figura y dos ejes. \texttt{ejes[0].plot(x,valores)} dibuja los pares $(x_j,\mathrm{valores}_j)$ en el panel izquierdo. \texttt{ejes[1]} identifica el derecho. La etiqueta \texttt{label} nombra la curva; \texttt{legend()} muestra esos nombres. \texttt{set\_xlabel} y \texttt{set\_ylabel} escriben qué significa cada eje. \texttt{grid} solo añade líneas de guía.

\paso{4. Qué se pasa a cada figura.}
\begin{center}\small\begin{tabular}{lll}\toprule
Archivo & Panel izquierdo & Panel derecho\\\midrule
\texttt{actividad.png} & $P_N$ & $e^{P_N}$\\
\texttt{errores.png} & $|P_N-F_3|$ & $100|e^{P_N-F_3}-1|$\\
\texttt{termodinamica.png} & $H_N$ & $g_N^E/1000$\\
\texttt{serie\_infinita.png} & $Q_m$ y $e^{F_3}$ & $100|Q_m/e^{F_3}-1|$\\\bottomrule
\end{tabular}\end{center}
Las figuras didácticas usan escalas lineales y muestran todos los órdenes, incluido el 0. Las figuras editoriales del PDF ajustan escalas y selección de curvas para destacar regiones de interés. Los cálculos son equivalentes; las presentaciones no son idénticas.

\paso{5. Guardar y ejecutar.}
\texttt{savefig} escribe la imagen; \texttt{dpi=180} indica resolución. \texttt{close} libera memoria. \texttt{main()} coordina cálculo, tablas y figuras. La condición \texttt{if \_\_name\_\_ == "\_\_main\_\_"} lo llama cuando se ejecuta el archivo; no lo llama al importarlo para pruebas. Cada instrucción del listado tiene un comentario inmediato que explica su finalidad.

\tema{e) Código didáctico completo, con comentarios instrucción por instrucción}
\textbf{Cómo copiarlo.} Preferir el archivo del repositorio, que conserva espacios, acentos y sangría. Los números del margen del listado son guías de lectura; no forman parte del código. Las líneas largas pueden partirse visualmente en el PDF sin representar instrucciones nuevas. El archivo completo se ejecuta con el comando indicado; no hay que agregar bloques que falten.
\lstinputlisting{ejemplos/paso_a_paso.py}

\tema{Versión MATLAB: archivo completo y forma de ejecución}
\paso{1. Qué archivo usar.}
La versión MATLAB completa se encuentra en \texttt{matlab/paso\_a\_paso.m}. Incluye una función principal y todas sus funciones auxiliares en el mismo archivo. No requiere copiar funciones desde otro archivo ni instalar un toolbox científico adicional. Los parámetros y las ecuaciones son los mismos que en Python.

\paso{2. Abrir la carpeta correcta.}
En MATLAB, seleccionar la raíz del repositorio como carpeta actual. En la ventana de comandos ejecutar:
\begin{verbatim}
addpath('matlab');
resultados = paso_a_paso;
\end{verbatim}
\texttt{addpath} permite que MATLAB encuentre el archivo dentro de la subcarpeta. La segunda línea llama a la función, realiza el cálculo completo y guarda en \texttt{resultados} una estructura con las magnitudes calculadas.

También se puede abrir el archivo y ejecutarlo desde su propia carpeta. El programa localiza su ruta mediante \texttt{mfilename('fullpath')}; la carpeta de salida no depende de dónde se abra la consola. No se deben pegar los comentarios del PDF como si fueran órdenes distintas: el listado ya está completo en el archivo descargable.

\paso{3. Resultados que se generan.}
La carpeta \texttt{output/matlab/} contiene cuatro imágenes. Sus nombres son \texttt{actividad.png}, \texttt{errores.png}, \texttt{termodinamica.png} y \texttt{serie\_infinita.png}. También se guardan \texttt{metricas.csv}, \texttt{coeficientes\_gamma.csv}, \texttt{taylor\_gamma.csv} y \texttt{malla.csv}. La malla contiene los 10001 valores de $x$ y todas las columnas de los polinomios, coeficientes, hierro, energía y expansión directa.

\paso{4. Consultar valores en la consola.}
\begin{verbatim}
resultados.metricas
resultados.b
resultados.P(4, end)
resultados.G(4, end)/1000
\end{verbatim}
La primera orden muestra las métricas. La segunda, $b_0$ hasta $b_{12}$. La tercera toma el orden 3 en la última composición; debe dar aproximadamente $-1.883613$. La cuarta devuelve la energía del mismo punto en kJ/mol, aproximadamente $-16.414661$. \texttt{end} significa la última posición de esa dimensión.

\paso{5. Elegir otra carpeta de salida.}
\texttt{resultados = paso\_a\_paso('mis\_resultados');} pasa un argumento: el nombre de la carpeta. Si no se da argumento, se usa la ruta predeterminada. Cambiar la carpeta no cambia parámetros ni cálculos.

\tema{MATLAB desde cero: las diferencias que pueden cambiar una cuenta}
\paso{1. Índices que comienzan en uno.}
Matemáticamente usamos $c_0,c_1,c_2,c_3$. En MATLAB se guardan en \texttt{C(1)}, \texttt{C(2)}, \texttt{C(3)}, \texttt{C(4)}. Por eso la instrucción para $c_k$ es \texttt{C(k+1)}. El contador matemático puede valer cero aunque las posiciones del arreglo no.

Para la recurrencia, $b_{n-k}$ se encuentra en \texttt{b(n-k+1)}. Ejemplo: con $n=2$, $k=2$, se necesita $b_0$ y la posición es $2-2+1=1$. El resultado $b_2$ se guarda en \texttt{b(3)}.

\paso{2. Potencias, productos y divisiones por componentes.}
\texttt{x.\^{}2} eleva cada entrada de $x$ al cuadrado. \texttt{x.*P} multiplica las entradas correspondientes de dos arreglos. \texttt{Q./referencia} divide cada entrada de $Q$ entre la correspondiente de la referencia. El punto importa: \texttt{*}, \texttt{/} y \texttt{\^{}} sin punto pueden significar operaciones matriciales diferentes. Multiplicar por un escalar, como \texttt{2.06*x.\^{}2}, sí usa \texttt{*} porque el factor es un solo número.

\paso{3. Filas y columnas.}
\texttt{zeros(4,M)} crea cuatro filas y $M$ columnas. \texttt{P(fila,:)} selecciona todas las columnas de una fila. Para orden 2, \texttt{fila=2+1=3}. \texttt{P(3,:)} representa $P_2$ en toda la malla. \texttt{P.'} intercambia filas y columnas sin conjugación compleja: se usa para guardar una composición por fila en CSV y para dibujar una curva por columna.

\paso{4. Ciclos y cierres.}
\texttt{for k=0:orden} recorre todos los enteros entre esos límites, incluidos ambos. \texttt{end} termina el bloque. Para orden 3, \texttt{(orden+1):3} es \texttt{4:3}, una secuencia vacía con incremento positivo predeterminado; la suma de términos omitidos queda en cero. El punto y coma al final de una instrucción evita imprimir automáticamente su resultado, pero no evita que se calcule.

\paso{5. Funciones y estructuras.}
\texttt{function p=logaritmo\_aluminio(x,orden,C)} define una función que recibe tres argumentos y devuelve \texttt{p}. \texttt{resultados.P} accede al campo \texttt{P} de una estructura. Las llaves de \texttt{\{'Orden 0','Orden 1'\}} crean una colección de textos. \texttt{...} al final de una línea indica que la instrucción continúa en la siguiente.

\tema{MATLAB: seguir las cuentas y comprobar la equivalencia}
\paso{1. Seguir el orden 2 en una composición.}
El ciclo empieza con $p=0$. Para $k=0$ usa \texttt{C(1)} y $x^0=1$; para $k=1$ usa \texttt{C(2)} y $x$; para $k=2$ usa \texttt{C(3)} y $x^2$. En $x=0.1$ acumula $-3.8685$, luego $-3.3912$ y finalmente $-3.3706$. Es la misma cuenta desarrollada antes.

\paso{2. Seguir la primera iteración de la recurrencia.}
\texttt{b(1)=exp(C(1))} fija $b_0$. Con $n=1$, el ciclo interior solo tiene $k=1$; calcula \texttt{1*C(2)*b(1)} y divide entre 1. Guarda el resultado en \texttt{b(2)}, que representa $b_1$. Con $n=2$ ya se pueden usar $b_1$ y $b_0$. Así nunca se necesita un coeficiente que aún no exista.

\paso{3. Conservar las mismas convenciones físicas.}
$G$ está en J/mol y solo se divide entre 1000 al dibujarlo en kJ/mol. $P$ y $H$ son logaritmos; \texttt{Gamma=exp(P)} obtiene el coeficiente. El error de la serie infinita usa \texttt{Q./referencia-1}, no otra exponencial. La cola del logaritmo se estabiliza con el mismo umbral $x<0.1$ y hasta grado 32.

\paso{4. Comparar automáticamente los archivos.}
Después de ejecutar la versión MATLAB o compatible, desde la raíz del repositorio:
\begin{verbatim}
uv run --locked python pruebas/verificar_matlab.py
\end{verbatim}
El verificador compara todas las columnas de los 10001 puntos, los 13 coeficientes de la recurrencia y ambas tablas de errores con el modelo Python. Permite diferencias pequeñas de redondeo entre motores. Si una magnitud no coincide dentro de las tolerancias indicadas en el verificador, termina con error.

\paso{5. Entorno de validación disponible.}
La ejecución de compatibilidad se realiza con GNU Octave, que interpreta este archivo de sintaxis MATLAB, usando \texttt{Dockerfile.octave}. No se dispone de una licencia ni de un motor MATLAB instalado en el entorno de preparación; por tanto la comprobación con Octave no se presenta como una ejecución nativa en MATLAB. El archivo utiliza funciones básicas y evita toolboxes adicionales.
\begin{verbatim}
docker build -f Dockerfile.octave -t rey-octave .
docker run --rm -v "$PWD:/proyecto" rey-octave
\end{verbatim}
En Linux se puede añadir \texttt{--user "\$(id -u):\$(id -g)"} a \texttt{docker run} para que los archivos pertenezcan al usuario actual. La receta fija la distribución base declarada, pero no archiva todas las versiones del sistema por digest. El README registra la versión del motor efectivamente comprobada.

\tema{Código MATLAB completo, comentado instrucción por instrucción}
Los comentarios comienzan con \texttt{\%}; explican la instrucción siguiente y no se ejecutan. Copiar preferentemente \texttt{matlab/paso\_a\_paso.m} desde el repositorio para conservar el archivo completo. Las funciones auxiliares están incluidas después del \texttt{end} de la principal.
\lstinputlisting[language=Matlab]{matlab/paso_a_paso.m}


\tema{Cómo funciona la reproducción editorial y cómo se comprueba}
\paso{1. Ubicación de cada responsabilidad.}
\texttt{src/feal/modelo.py} contiene las funciones matemáticas reutilizables. \texttt{reproducir.py} evalúa la malla, escribe tablas, produce figuras y llama a LaTeX. \texttt{desglose.py} escribe los pasos intermedios y los CSV de auditoría. El ejemplo didáctico hace explícitos los ciclos y acumulaciones que el módulo compacto expresa con sumas o funciones de NumPy.

\paso{2. Equivalencia con la evaluación compacta.}
El módulo reutilizable usa \texttt{np.polynomial.polynomial.polyval} para evaluar un polinomio. Por ejemplo,
\[
c_0+c_1x+c_2x^2+c_3x^3
=c_0+x[c_1+x(c_2+xc_3)].
\]
Al distribuir los paréntesis interiores aparecen $c_2x+c_3x^2$; al multiplicar nuevamente por $x$ y sumar $c_1x$ aparecen $c_1x+c_2x^2+c_3x^3$. La forma anidada reduce operaciones. Puede diferir en los últimos bits por el orden de redondeo, pero representa el mismo polinomio que el ciclo didáctico.

\paso{3. Tablas y texto sincronizados.}
Los comandos \texttt{input} de LaTeX insertan las tablas generadas, incluidas las multiplicaciones de la recurrencia. Las figuras se leen de los PDF vectoriales calculados. LaTeX se ejecuta tres veces para resolver el índice y las referencias a páginas. El PDF final se escribe en \texttt{output/pdf/solucionario.pdf}.

\paso{4. Verificar en vez de confiar solo en la apariencia.}
\begin{verbatim}
uv run --locked pytest
\end{verbatim}
Las pruebas verifican las derivadas mediante álgebra simbólica exacta, Gibbs-Duhem para cada orden, el hierro frente a integración numérica independiente, las métricas de referencia y la recurrencia frente a una expansión simbólica de la exponencial. También comparan el programa didáctico con el módulo reutilizable y verifican la estructura y operaciones de los archivos de desglose.

\paso{5. Reproducibilidad.}
\texttt{uv.lock} fija bibliotecas y \texttt{.python-version} el intérprete. Se fijan las fechas internas del PDF y se registran huellas SHA-256 en \texttt{output/manifiesto.json}. Las salidas del ejemplo didáctico se guardan aparte y no forman parte de ese manifiesto editorial. Las versiones de TeX Live y las fuentes dependen del sistema: no se promete identidad binaria entre plataformas diferentes. En el mismo entorno se compara una reconstrucción limpia y una segunda ejecución.

\tema{Guía para responder las dudas que motivaron esta edición}
\paso{«¿Dónde está Taylor?»}
En la construcción $c_k=F^{(k)}(0)/k!$, con las derivadas calculadas término a término y evaluadas en cero. La parte previa deriva la función que se va a desarrollar; esa parte previa no debe confundirse con Taylor.

\paso{«¿Por qué sale lo mismo que ya estaba?»}
Porque el modelo del logaritmo es un polinomio cúbico. Su Taylor de grado tres lo recupera. Cortar antes sí crea una aproximación con error calculable.

\paso{«¿De dónde toma los datos para las gráficas?»}
De evaluar las ecuaciones en una malla que elegimos. Solo los parámetros de partida proceden del artículo; los miles de puntos no son miles de observaciones experimentales.

\paso{«¿De dónde toma los coeficientes de la serie infinita?»}
De $b_0=e^{c_0}$ y la recurrencia que se deduce al igualar potencias en $\gamma'=F_3'\gamma$. Se muestran todos los coeficientes utilizados, con sus productos, sumas y divisiones.

\paso{«¿El fragmento corto era todo el programa?»}
No. El fragmento del PDF original era solo un núcleo. Esta edición incorpora un archivo ejecutable completo y su listado comentado; incluye las instrucciones para generar y guardar las gráficas.

\paso{«¿Puedo reproducir un número específico?»}
Sí: identificar si es $P_N$, $\gamma_N$, $H_N$, $g_N^E$ o $Q_m$; localizar su definición; sustituir el $x$ correspondiente; seguir el ejemplo y consultar el CSV. Para $P_N$ y sus errores, \texttt{desglose\_malla.csv} deja explícitos los términos y sus combinaciones. Para la serie infinita, consultar \texttt{coeficientes\_gamma.csv} y \texttt{terminos\_gamma\_extremo.csv}.

\tema{Fuentes y correspondencia con el examen}
[1] Ishfaq, M.; Pande, M. M. (2025). \textit{Thermodynamic analysis of Fe-Al-O system using Redlich-Kister (R-K) polynomial and Wagner's Taylor series expansion}. The Journal of Chemical Thermodynamics, 202, 107415. DOI: \href{https://doi.org/10.1016/j.jct.2024.107415}{10.1016/j.jct.2024.107415}. Copia original: \texttt{referencias/articulo.pdf}. Parámetros: tabla 1, p. 5. Modelo y relaciones: ecuaciones (16)--(26), p. 6. Composiciones de referencia: tabla 2, p. 7. Reducción binaria: sección 3.2.

[2] Universidad Veracruzana, Facultad de Ciencias Químicas. \textit{Primer Examen Parcial: Matemáticas Aplicadas}. Maestría en Ingeniería Química, 6 de octubre de 2026. Copia: \texttt{referencias/parcial.pdf}.

[3] \textit{Solucionario razonado: Series de Taylor y termodinámica del sistema Fe-Al-O}. Documento proporcionado, fechado el 7 de octubre de 2026. Se conserva sin cambios en la carpeta \texttt{referencias/}, con el nombre \texttt{solucionario\_original.pdf}. Esta edición amplía sus explicaciones, reconstruye sus resultados y agrega código íntegro y desgloses; no reproduce su maquetación original.

\begin{center}\begin{tabular}{lp{0.76\linewidth}}\toprule
Inciso & Contenido desarrollado\\\midrule
a & Fenómeno de desoxidación, variables, actividades, unidades y aproximación binaria.\\
b & Elección de Taylor, función a desarrollar, centro en cero y diferencia con Fourier.\\
c & Derivación parcial molar, álgebra del modelo, derivadas de Taylor, factoriales y Gibbs-Duhem.\\
d & Truncamientos, ejemplos, errores, máximos, RMSE y tolerancias.\\
e & Programa completo, instrucciones de ejecución, gráficas, código comentado y pruebas.\\
f & Interpretación de tendencias, energía, estabilidad y límites del modelo.\\
g & Conclusiones fundamentadas, condiciones de uso y límites de las afirmaciones.\\\bottomrule
\end{tabular}\end{center}

\textbf{Convención final.} Se usa punto decimal para coincidir con Python y con las cifras de las ecuaciones. Los valores del texto se redondean solo al imprimir. Los originales se mantienen íntegros y sus huellas constan en \texttt{referencias/SHA256SUMS}.
